\subsection{Construction order as a design variable}\label{sec:order}

The path dependence above has a practical reading. If the same
final genotype can be reached by different orders of introduction,
and the orders do not commute, then which order to choose is a
design decision --- in strain construction, in directed evolution,
in any protocol that builds a genotype in stages. The
double-mutant panels quantify it. Among non-synthetic-lethal
pairs, the 90th-percentile endpoint separation between the two
construction orders is $101.3$ in $\ell_1$ flux units, against a
wild-type total flux magnitude of $770$: the same final genotype
can sit roughly $13\%$ of total flux magnitude away from the other
order's state.

Is the better order predictable in advance? Yes --- from a
quantity the singles panel already computes. For each of the $66$
non-degenerate closed loops, let $P$ be the normalized difference
in single-knockout curvature between the two genes (which knockout
crosses more wall mass), and $A$ the normalized asymmetry in
terminal drift between the two excursion orders. They track each
other: Spearman $\rho = 0.442$ ($p = 2.0 \times 10^{-4}$), with
sign agreement in $59$ of $66$ pairs (binomial $p = 2.4 \times
10^{-11}$). The gene whose knockout crosses more wall mass sets
the loop's irreversibility when it leads the excursion. The
planning rule is one line: schedule deep wall-crossers last to
minimize terminal drift --- the right choice when intermediate
states must remain recoverable --- or first to lock the endpoint
away from the initial state, when a production phenotype should
persist.

An honest null sharpens what the rule is not. Order asymmetry is
not predicted by active-set footprint overlap ($\rho = 0.07$,
$p = 0.38$, against the same Jaccard index that tracks epistasis
magnitude at $\rho_S = 0.865$): epistasis alignment and order
ranking are governed by different geometric quantities ---
overlap versus mass --- and the machinery for either is a single
LP per candidate gene, already a byproduct of the singles panel.
